\section{数据主权与 Data Embassy：位置、控制和恢复是不同问题}

国家级系统会遇到 data sovereignty。课堂提到 data embassy 和虚拟化，但这类概念很容易被简化成“数据放在国内/国外”。本节只讨论架构维度，不从访谈推导具体法律结论：location、control、jurisdiction、portability 与 continuity 必须分别设计，并通过合同、技术控制和演练共同验证。

\subsection{五个问题替代一个 location checkbox}

\term{data sovereignty}（数据主权）描述数据受哪些国家法律、控制权与政策约束；\term{data residency}（数据驻留）更窄，关注数据物理或逻辑存放位置；\term{data embassy}（数据使馆）通常指通过特殊法律安排和跨境基础设施，为关键数字资产提供主权保护与连续性。它不是普遍标准产品，必须阅读具体条约和法律文本。

\lecturefigure{13-data-sovereignty.png}{数据主权需要同时设计 location、control、jurisdiction 与 continuity。}{本地字幕 32:20--36:50；Saudi DGA cloud-governance 资料背景。}

\begin{knowledgebox}{读图：架构评审的五个问题}
数据在哪里静态存储和处理？谁持有 keys 与 privileged identity？哪个司法辖区可以提出 disclosure？数据和模型能否 portable export？在 provider、region 或国家级网络故障下如何恢复？只有五个问题都有明确 owner 和 test，sovereignty 才从口号变成控制系统。
\end{knowledgebox}

\begin{warningbox}{位置不会自动带来安全或可用性}
单一区域本地化可能增加 correlated failure；多区域复制又可能引入跨境和一致性问题。Encryption、key custody、identity、logging、supply-chain assurance、backup integrity 与 restore drill 必须与 location policy 一起设计。
\end{warningbox}

\teachervoice{课堂的价值在于把国家韧性视为 infrastructure requirement，而不是事后合规。讲义进一步补充：任何“数字使馆”方案都必须把法律承诺翻译成 key ownership、replication topology、failover authority 和定期演练。}

\subsection{本章小结}

Data sovereignty 是多维架构问题。Residency 只是其中一个输入；控制权、司法辖区、portability、恢复和审计决定系统在危机时是否真正可用。

